\chapter{Indices, denominators and elliptic coefficient lines}
\label{cyidx:chapter}

Taking an Euler characteristic turns a complex of geometric
coefficients into a number. For Hilbert schemes of a K3 surface
these numbers assemble into the reciprocal of an Igusa product,
after a boundary factor is removed. An Enriques surface gives
the corresponding half-index. The root algebra supplies another
source of the same product, with a Lie bracket and a Chevalley
differential. Comparing the two constructions requires the
coefficient complexes and their operations.

The boundary factor is already significant at length zero.
The elliptic coefficient line is significant at a collision:
its section can have a double zero, and the resulting quotient
retains a normal nilpotent. For three interacting colours,
the simultaneous incidence has a nonzero higher multiplication.
We will compute this complex, its deformation over the spectral
elliptic curve, and the cover on which its coefficient line
has a square root. These constructions retain data that their
Euler characteristics forget.

\section{The elliptic index and its coefficient complexes}
\label{cyidx:sec:index}

Let \(Y\) be a smooth projective complex variety of dimension \(d\).
For a vector bundle \(V\), write
\[
 \Lambda_zV=\sum_{j\geq0}z^j\Lambda^jV,\qquad
 \operatorname{Sym}_zV=\sum_{j\geq0}z^j\operatorname{Sym}^jV.
\]
The elliptic class is the formal virtual bundle
\begin{equation}\label{cyidx:eq:elliptic-class}
\begin{split}
 \mathcal E_{q,y}(Y)=y^{-d/2}\bigotimes_{r\geq1}
 \bigl(&\Lambda_{-yq^{r-1}}T_Y^*
       \otimes\Lambda_{-y^{-1}q^r}T_Y\\
       &\otimes\operatorname{Sym}_{q^r}T_Y^*
       \otimes\operatorname{Sym}_{q^r}T_Y\bigr)
 \in K^0(Y)[y^{\pm1/2}][[q]].
\end{split}
\end{equation}
At a fixed power of \(q\), only finitely many positive-energy
factors contribute. The remaining exterior factor has finite
rank, so the Laurent support in \(y\) is finite as well.
Define
\[
 \operatorname{Ell}(Y;q,y)=
 \chi(Y,\mathcal E_{q,y}(Y)),\qquad
 \chi(Y,V)=\sum_j(-1)^j\dim H^j(Y,V),
\]
coefficient by coefficient.

There is a complex before this Euler characteristic. Give each
exterior power in \eqref{cyidx:eq:elliptic-class} its exterior
degree modulo two, keep the symmetric powers even, and retain
the displayed \(q,y\) weights. This defines a formal superbundle
\(\mathscr E_{q,y}(Y)\) whose super-\(K\)-class is
\(\mathcal E_{q,y}(Y)\). Its coefficient \(\mathscr E_{m,\ell}(Y)\)
is a finite sum of bundles with specified internal parity.
The complex
\begin{equation}\label{cyidx:eq:coefficients}
 \mathscr V_{m,\ell}(Y)=
 R\Gamma(Y,\mathscr E_{m,\ell}(Y))
\end{equation}
is perfect over \(\mathbb C\). Its Euler superdimension uses
\((-1)^{j+\epsilon}\), where \(j\) is cohomological degree and
\(\epsilon\) is internal parity.

Fix a K3 surface \(S\) and an Enriques surface \(T\).
Write \(\widetilde T\to T\) for the canonical étale double
cover. The K3 elliptic genus is
\begin{equation}\label{cyidx:eq:seed}
 \operatorname{Ell}(S;q,y)=2\phi_{0,1}(q,y),\qquad
 \phi_{0,1}(\tau,z)=
 4\sum_{a=2}^4\frac{\vartheta_a(\tau,z)^2}
                         {\vartheta_a(\tau,0)^2},
\end{equation}
with \(q=e^{2\pi i\tau}\), \(y=e^{2\pi iz}\).
The theta normalization is that of \cite{EZ}.
Write
\[
 \phi_{0,1}(q,y)=\sum_{m\geq0,\ell\in\mathbb Z}
                      f(4m-\ell^2)q^my^\ell.
\]
The elliptic transformation law identifies the coefficients
at \((m,\ell)\) and \((m+a\ell+a^2,\ell+2a)\).
Reducing \(\ell\) to zero or one proves that \(f(D)=0\)
for \(D<-1\). The constant row is
\[
 [q^0]\phi_{0,1}=y^{-1}+10+y.
\]
Indeed the K3 Hodge numbers give
\(y^{-1}(\chi(\mathcal O_S)-y\chi(\Omega_S^1)
+y^2\chi(\Omega_S^2))=2y^{-1}+20+2y\).
In particular \(f(-1)=1\), \(f(0)=10\), whereas
\(\phi_{0,1}(\tau,0)=12\).

\begin{proposition}\label{cyidx:prop:half-index}
The Enriques elliptic genus is \(\phi_{0,1}\).
At \(y=1\), the elliptic genus of any smooth projective
\(Y\) is its topological Euler characteristic, independent
of \(q\).
\end{proposition}
\begin{proof}
The tangent bundle pulls back along an étale map. The
Hirzebruch--Riemann--Roch characteristic class defining
the elliptic genus therefore pulls back, and its integral
over the double cover is twice its integral over \(T\).
Apply \eqref{cyidx:eq:seed} to \(\widetilde T\).

At \(y=1\), use
\(\Lambda_{-z}V\,\operatorname{Sym}_zV=1\).
All positive-\(q\) factors cancel, leaving
\(\Lambda_{-1}T_Y^*\). Its holomorphic Euler characteristic
is \(\sum_{p,j}(-1)^{p+j}h^{p,j}(Y)=e(Y)\).
\end{proof}

\subsection{Length and the boundary factor}

For a surface \(Y\), let \(Y^{[N]}\) be its Hilbert scheme
of length-\(N\) subschemes, and put
\[
 F_Y(p,q,y)=\sum_{N\geq0}p^N
                  \operatorname{Ell}(Y^{[N]};q,y).
\]
Thus \(p\) records length, \(q\) elliptic energy, and
\(y\) the Laurent weight. Each coefficient has the
complex \eqref{cyidx:eq:coefficients} on \(Y^{[N]}\);
the empty Hilbert scheme contributes \(1\).

For triples \((a,b,\ell)\), the notation \((a,b,\ell)>0\)
means \(a,b\geq0\), \(a+b>0\), or \(a=b=0\), \(\ell<0\).
Choose square roots of the formal monomials and define
\begin{equation}\label{cyidx:eq:D5}
 D_5=(pqy)^{1/2}
       \prod_{(a,b,\ell)>0}
          (1-p^aq^by^\ell)^{f(4ab-\ell^2)},\qquad
 \Xi=D_5^2.
\end{equation}
The length-zero factor, including its Weyl monomial, is
\begin{equation}\label{cyidx:eq:boundary}
\begin{split}
 b(q,y)&=q^{1/2}(y^{1/2}-y^{-1/2})
       \prod_{r\geq1}(1-q^r)^{10}
                         (1-yq^r)(1-y^{-1}q^r),\\
 \mathcal B(q,y)&=b(q,y)^2\\
 &=q(y+y^{-1}-2)\prod_{r\geq1}
        (1-q^r)^{20}(1-yq^r)^2(1-y^{-1}q^r)^2.
\end{split}
\end{equation}
Here the symbol \(b(q,y)\) is a function; the integer \(b\)
in a root triple is its energy coordinate.

\begin{theorem}\label{cyidx:thm:Hilbert-product}
In \(\mathbb Z[y,y^{-1}][[p,q]]\),
\begin{align}
 F_T&=\prod_{a>0,b\geq0,\ell\in\mathbb Z}
       (1-p^aq^by^\ell)^{-f(4ab-\ell^2)}
       =\frac{p^{1/2}b(q,y)}{D_5},\label{cyidx:eq:Enriques}\\
 F_S&=\prod_{a>0,b\geq0,\ell\in\mathbb Z}
       (1-p^aq^by^\ell)^{-2f(4ab-\ell^2)}
       =\frac{p\mathcal B(q,y)}{\Xi}=F_T^2.
       \label{cyidx:eq:K3}
\end{align}
The quotients are expanded after canceling their
length-zero factors.
\end{theorem}
\begin{proof}
The smooth-surface Hilbert-scheme formula of
Borisov--Libgober identifies the elliptic genus of the
Hilbert--Chow resolution with the symmetric-product
orbifold genus \cite{rewriteBL2005}*{Corollary~6.7}.
It applies to both surfaces; triviality of the canonical
bundle is not a hypothesis of that formula. Their seed
coefficients are \(f\) and \(2f\), respectively.
At fixed \(p^Nq^M\), only \(a\leq N\), \(b\leq M\)
contribute, and \(4ab-\ell^2\geq-1\) bounds \(\ell\).
This proves the formal finiteness of the products.
Extracting the factors with \(a=0\) from
\eqref{cyidx:eq:D5} gives \eqref{cyidx:eq:boundary}.
Squaring doubles the positive-length exponents.
\end{proof}

The first coefficients make the normalization visible:
\begin{align}
 F_T(p,q,1)&=1+12p+90p^2+520p^3+O(p^4),\label{cyidx:eq:Euler-T}\\
 F_S(p,q,1)&=1+24p+324p^2+3200p^3+O(p^4).
 \label{cyidx:eq:Euler-S}
\end{align}
They are the products
\(\prod_{a\geq1}(1-p^a)^{-12}\) and
\(\prod_{a\geq1}(1-p^a)^{-24}\).
For a more refined coefficient, put \(A(y)=y+10+y^{-1}\).
The length-one modes contribute
\(\tfrac12(A(y)^2+A(y^2))\) at length two, and the
length-two modes contribute \(A(y)\). Hence
\begin{align}
 [p^2q^0]F_T&=y^2+11y+66+11y^{-1}+y^{-2},\label{cyidx:eq:two-T}\\
 [p^2q^0]F_S&=3y^2+42y+234+42y^{-1}+3y^{-2}.
 \label{cyidx:eq:two-S}
\end{align}
Both \(\mathcal B\) and \(\Xi\) vanish at \(y=1\).
Their canceled quotient has the regular specialization
\eqref{cyidx:eq:Euler-S}; their separate reciprocals do not.

The complexes \(\mathscr V_{m,\ell}(T^{[N]})\) provide
geometric coefficients for \(F_T\). Formula
\eqref{cyidx:eq:K3} compares their Euler supercharacters.
It supplies neither a tensor-square equivalence with
the K3 complexes nor extension pull--push maps on them.
A perfect complex can have zero Euler characteristic
and nonzero cohomology, so taking the index cannot
recover such maps.

\section{The positive-length root ideal}
\label{cyidx:sec:root-ideal}

Use the lattice \(L=\mathbb Z^3\) with coordinates
\((a,b,\ell)\) and pairing
\[
 ((a,b,\ell),(a',b',\ell'))=
 2\ell\ell'-4(ab'+ba').
\]
Its real simple roots and Weyl vector are
\begin{equation}\label{cyidx:eq:real-roots}
 \delta_1=(0,1,1),\quad
 \delta_2=(1,0,1),\quad
 \delta_3=(0,0,-1),\qquad
 \rho=(1/2,1/2,1/2).
\end{equation}
Their Gram matrix is \(2\) on the diagonal and \(-2\)
off the diagonal, and \((\rho,\delta_i)=-1\).
Fix
\[
 e^{-\delta_1}=qy,\qquad e^{-\delta_2}=py,\qquad
 e^{-\delta_3}=y^{-1},\qquad e^{-\rho}=(pqy)^{1/2}.
\]
These are root-torus coordinates; the analytic doubling
in the Gritsenko--Nikulin convention will be specified below.

Let \(\mathfrak n_+\) be the positive part of the
Gritsenko--Nikulin Lie superalgebra, with the full labelled
imaginary simple-root spaces and parities of its presentation
\cite{GN97}*{Sections~3 and~6}.
The real generators are even. The presentation supplies
both parity spaces and the bracket. The denominator theorem
gives their difference:
\[
 \operatorname{sdim}(\mathfrak n_+)_{a,b,\ell}
       =f(4ab-\ell^2).
\]
Here \(\operatorname{sdim}V=\dim V_{\bar0}-\dim V_{\bar1}\).
Actual root spaces are finite-dimensional in each parity.

Length \(a\) defines
\[
 \mathfrak u=\bigoplus_{a>0,b,\ell}(\mathfrak n_+)_{a,b,\ell},
 \qquad
 \mathfrak n_0=\bigoplus_{b,\ell}(\mathfrak n_+)_{0,b,\ell}.
\]
For Chevalley chains, internal parity and chain degree
are separate. Write
\[
 \chi_s C_*(\mathfrak u)
       =\sum_{j\geq0}(-1)^j\operatorname{sch}C_j(\mathfrak u).
\]
All characters are formed degree by degree in the positive
root cone.

\begin{proposition}\label{cyidx:prop:PBW}
There is a semidirect decomposition
\(\mathfrak n_+=\mathfrak u\rtimes\mathfrak n_0\), and
\begin{equation}\label{cyidx:eq:PBW}
 \operatorname{sch}U(\mathfrak u)=F_T,\qquad
 \chi_s C_*(\mathfrak u)=F_T^{-1},\qquad
 D_5=p^{1/2}b(q,y)\,\chi_s C_*(\mathfrak u).
\end{equation}
\end{proposition}
\begin{proof}
The bracket adds the nonnegative length. Its positive
part is therefore an ideal, and its zero part a subalgebra.
Their vector-space sum is direct.

In super PBW, an even root vector of monomial weight \(x\)
contributes \((1-x)^{-1}\), and an odd root vector contributes
\(1-x\). In Chevalley chains the contributions are inverse.
For an odd vector, symmetric powers have equal internal
and chain parities, whose signs cancel; for an even vector,
the chain sign gives its exterior factor \(1-x\).
The root supermultiplicities and
\eqref{cyidx:eq:Enriques} prove the first two identities.
Extracting the length-zero part gives the third.
\end{proof}

The full enveloping character is \((pqy)^{1/2}/D_5\);
the character of \(U(\mathfrak n_0)\) is
\((qy)^{1/2}/b(q,y)\). Their quotient is \(F_T\).
The ideal still carries the adjoint action of
\(\mathfrak n_0\). Removing a scalar factor from a character
does not discard that action from the algebra.

\subsection{Completion by the total height of a monomial}

Every root charge has the expression
\begin{equation}\label{cyidx:eq:height}
 (a,b,\ell)=b\delta_1+a\delta_2+(a+b-\ell)\delta_3,
 \qquad h(a,b,\ell)=2(a+b)-\ell.
\end{equation}
On the positive cone this is a positive integral height.
For an active product factor with \(a+b>0\), the support
bound implies \(\ell\leq a+b\), hence \(h\geq a+b\).
Indeed \(\ell\geq a+b+1\) would give
\(\ell^2-4ab\geq(a-b)^2+2(a+b)+1>1\).
The only active factor with \(a=b=0\) has \(\ell=-1\).
Bounded height therefore contains finitely many root degrees.

For a height-graded vector space \(V=\bigoplus_{r\geq0}V_r\),
set
\[
 \widehat V=\prod_{r\geq0}V_r,\qquad
 \widehat V\,\widehat\otimes\,\widehat W
 =\prod_{r\geq0}\bigoplus_{i+j=r}V_i\otimes W_j.
\]
The neighborhoods of zero are the products over \(r\geq H\).
The height of a monomial in \(U(\mathfrak n_+)\) is the
sum of the heights of all its generators, including repetitions.
The reduced bar and Chevalley complexes inherit this grading.

\begin{theorem}\label{cyidx:thm:completed-bar}
Antisymmetrization gives a quasi-isomorphism
\begin{equation}\label{cyidx:eq:completed-bar}
 \prod_{r\geq0}C_*(\mathfrak n_+)_r
 \longrightarrow
 \prod_{r\geq0}B_*(U(\mathfrak n_+))_r.
\end{equation}
At each total height, the terms are finite-dimensional and
only finitely many chain degrees occur.
\end{theorem}
\begin{proof}
A monomial of height \(r\) contains at most \(r\)
positive-height generators. Only finitely many root spaces
can occur, and each is finite-dimensional in both parities.
This proves both finiteness assertions, including the bar
length bound.

Use the bar differential whose first multiplication term
is positive and the Chevalley differential sending
\(x\wedge y\) to \([x,y]\). Antisymmetrization sums the
bar words with the permutation sign and the Koszul sign
of the internal parities. Adjacent multiplication terms
cancel in pairs except for the brackets, so it is a chain
map. Its PBW associated graded is the Koszul comparison
for a free supercommutative algebra: exterior chains for
even generators, symmetric chains for odd generators.
The filtration is finite at each total height.
Thus the comparison is a quasi-isomorphism there.
Products of complexes of complex vector spaces are exact,
so the product of these comparisons is again a
quasi-isomorphism.
\end{proof}

Truncating the Lie algebra alone gives a different topology.
For \(\mathfrak a=\mathbb Cx\), with \(x\) even and \(h(x)=1\),
\begin{equation}\label{cyidx:eq:completion-test}
 \varprojlim_H U\!\left(\mathfrak a/
                       \bigoplus_{r>H}\mathfrak a_r\right)
       =\mathbb C[x],\qquad
 \widehat{U(\mathfrak a)}=\mathbb C[[x]].
\end{equation}
The Lie quotient is constant for \(H\geq1\).
Total-height completion also permits infinite powers
of this one fixed generator.

For the level-zero current envelope, the same distinction
applies to products of negative modes. Total root height
defines a pro-vertex algebra by its quotients; at a fixed
height its states remain algebraic in the mode variables.
The coefficient of a fixed height and conformal energy
is finite-dimensional. Uniform locality for arbitrary
infinite sums is a further requirement if one seeks a
vertex algebra on the completed vector space.
The ordinary comparison \eqref{cyidx:eq:completed-bar}
does not construct that locality or the supported maps
on collision configurations.

\section{The automorphic section and its roots}
\label{cyidx:sec:automorphic-line}

Let
\[
 \mathfrak H_2=
 \left\{Z=\begin{pmatrix}\tau&z\\z&\sigma\end{pmatrix}
                  :\operatorname{Im}Z>0\right\}.
\]
Use \(q=e^{2\pi i\tau}\), \(y=e^{2\pi iz}\),
\(p=e^{2\pi i\sigma}\). For even characteristics
\(a,b\in\{0,1\}^2\), meaning \(a^{\mathsf T}b\) is even,
fix the Gritsenko--Nikulin phase convention
\begin{equation}\label{cyidx:eq:theta-phase}
 \vartheta^{\mathrm{GN}}_{a,b}(Z)=
 \sum_{m\in\mathbb Z^2}
  \exp\!\left(\pi i\left[
    (m+a/2)^{\mathsf T}Z(m+a/2)+b^{\mathsf T}m\right]\right).
\end{equation}
There are ten such characteristics.
This formula fixes the phases as well as the characteristics.

\begin{theorem}\label{cyidx:thm:automorphic}
The product \(D_5\) of \eqref{cyidx:eq:D5} converges
on a cusp chamber and extends to a holomorphic
Siegel modular form of weight \(5\), with a character
\(\nu_5\) of order two. In the convention
\eqref{cyidx:eq:theta-phase},
\begin{equation}\label{cyidx:eq:normalization}
 \Delta_5=\prod_{a^{\mathsf T}b\equiv0\ (2)}
                \vartheta^{\mathrm{GN}}_{a,b}=64D_5,
 \qquad \Xi=D_5^2=2^{-12}\Delta_5^2.
\end{equation}
The divisor of \(D_5\) in \(\mathfrak H_2\) is the
\(\operatorname{Sp}_4(\mathbb Z)\)-orbit of \(z=0\),
with multiplicity one.
\end{theorem}
\begin{proof}
The product, weight, character and divisor are the
Borcherds--Gritsenko--Nikulin theorem
\cite{GN97}*{Proposition~3.1 and Theorem~4.1}.
The weight is \(f(0)/2=5\); the negative-discriminant
coefficient \(f(-1)=1\) supplies the diagonal divisor.

The phase convention is \cite{GN97}*{equation~(1.2)}.
The four characteristics with \(a=0\) have leading
coefficient \(1\). The two with \(a=(1,0)\) contribute
\(2q^{1/8}\) each, and those with \(a=(0,1)\)
contribute \(2p^{1/8}\) each.
For \(a=(1,1)\), their lowest \(q,p\) terms are
\[
 2q^{1/8}p^{1/8}(y^{1/4}+y^{-1/4}),\qquad
 2q^{1/8}p^{1/8}(y^{1/4}-y^{-1/4}).
\]
Their product gives the leading Laurent factor
\(64q^{1/2}p^{1/2}(y^{1/2}-y^{-1/2})\).
This is \(64\) times the corresponding factor of
\eqref{cyidx:eq:D5}, fixing the scalar.
Squaring removes the quadratic character.
\end{proof}

The formal root variables in \eqref{cyidx:eq:real-roots}
give \(D_5\) directly. In the Lorentzian analytic
coordinates of the same Gritsenko--Nikulin construction,
the root denominator is \(D_5(2Z)\).
The map \(Z\mapsto Z/2\) relates these coordinate
conventions. It does not alter the root multiplicities.

The first Fourier--Jacobi coefficient is \(b(q,y)\).
For the odd theta function normalized by
\[
 \vartheta(\tau,z)=q^{1/8}(y^{1/2}-y^{-1/2})
    \prod_{r\geq1}(1-q^r)(1-yq^r)(1-y^{-1}q^r),
\]
it is \(b=\eta(\tau)^9\vartheta(\tau,z)\).
Its weight is \(5\) and its index is \(1/2\).
The weight-zero, index-one form \(\phi_{0,1}\)
is the multiplicative input; this first coefficient
is the additive-lift input. Their roles and weights differ.

\subsection{A line, a section and a trace}

On the moduli stack
\(\mathcal A_2=[\mathfrak H_2/\operatorname{Sp}_4(\mathbb Z)]\),
let \(\lambda\) be the determinant of the Hodge bundle,
with the convention that weight-\(k\) forms are sections
of \(\lambda^k\). The character \(\nu_5\) also denotes
its flat line. Thus
\begin{equation}\label{cyidx:eq:automorphic-line}
 D_5\in H^0(\mathcal A_2,\lambda^5\otimes\nu_5),
 \qquad
 \Xi\in H^0(\mathcal A_2,\lambda^{10}).
\end{equation}
The square relation identifies
\((\lambda^5\otimes\nu_5)^{\otimes2}\) with
\(\lambda^{10}\). The full line retains its Hodge
factor; its character alone is only the flat factor.

Roots of functions and roots of lines have different
descent data. If a nowhere-zero function \(f\) has local
\(n\)-th roots \(s_i\), their ratios \(s_i/s_j\)
form a \(\mu_n\)-valued one-cocycle. It defines the
torsor of roots in \(H^1(-,\mu_n)\).
For a line with transition functions \(g_{ij}\),
choose local \(n\)-th roots \(h_{ij}\).
The products \(h_{ij}h_{jk}h_{ki}\) form a
\(\mu_n\)-valued two-cocycle, the Kummer obstruction
in \(H^2(-,\mu_n)\) to an \(n\)-th root line.
The local functions \(s_i\) are generally
\(\mathbb G_m\)-valued, so they do not by themselves
trivialize the first cocycle in the sheaf \(\mu_n\).

Near a smooth point of the diagonal divisor, choose
a transverse coordinate \(u\) and trivialize the line.
Then \(D_5=u\,v(u)\), with \(v(0)\ne0\), and
\(\Xi=u^2v(u)^2\). A branch of \(\Xi^{1/n}\)
on the punctured disc has monodromy \(e^{4\pi i/n}\),
of order \(n/\gcd(n,2)\). In particular an eighth
root of \(\Xi\) has local order four, while an
eighth root of \(D_5\) has local order eight.
These local branching orders do not specify an
orientation line on compact sheaf moduli.

\begin{proposition}\label{cyidx:prop:trace}
Let \(V\) be a bounded finite-dimensional complex and
\(s:V\to V\) a chain involution. Its supertrace is an
integer. A homotopy equivalence intertwining two such
endomorphisms up to chain homotopy preserves their
supertraces. An arbitrary intertwining map need not do so.
\end{proposition}
\begin{proof}
The supertrace is the alternating trace on cohomology.
The eigenvalues of an involution in characteristic zero
are \(1\) and \(-1\), proving integrality.
A homotopy equivalence induces an isomorphism on cohomology,
and the intertwining homotopy makes the induced
endomorphisms conjugate. Their traces agree.
The inclusion \(\mathbb C\hookrightarrow\mathbb C^2\)
intertwines the identity endomorphisms, whose traces
are \(1\) and \(2\).
\end{proof}

Consequently a half-integral automorphic weight cannot
be the ordinary supertrace of an involution on a perfect
complex over \(\mathbb C\). An infinite graded trace,
a regularization or another coefficient-valued trace
requires its own definition. A duality from one object
to a different object also requires an identification
before it becomes an endomorphism to which a trace applies.

\section{Products at higher level}
\label{cyidx:sec:families}

At full level the weight is half the constant coefficient
of one Jacobi form. At a congruence level every cusp can
contribute. This distinction explains the weights of the
forms with a simple diagonal divisor.

For positive integers \(t,N\), set
\[
 \Gamma_t(N)=\operatorname{Sp}_4(\mathbb Q)\cap
 \begin{pmatrix}
 \mathbb Z&t\mathbb Z&\mathbb Z&\mathbb Z\\
 \mathbb Z&\mathbb Z&\mathbb Z&t^{-1}\mathbb Z\\
 N\mathbb Z&Nt\mathbb Z&\mathbb Z&\mathbb Z\\
 Nt\mathbb Z&Nt\mathbb Z&t\mathbb Z&\mathbb Z
 \end{pmatrix}.
\]
Its extension \(\Gamma_t(N)^+\) adjoins the involution
\((\tau,z,\sigma)\mapsto(t\sigma,z,\tau/t)\).
For \(t=1\), the first group is the Siegel congruence
subgroup with lower-left block divisible by \(N\).

Let \(\phi\) be a nearly holomorphic Jacobi form of
weight zero and index \(t\) on \(\Gamma_0(N)\).
Choose a representative \(f/e\) of each cusp, with
\(e\mid N\) and \((f,e)=1\), and a matrix taking
\(\infty\) to that cusp. Write its transformed expansion as
\[
 \phi_{f/e}(\tau,z)=\sum_{n,\ell}c_{f/e}(n,\ell)q^ny^\ell,
 \qquad h_e=\frac{N}{\gcd(e^2,N)},\qquad N_e=\frac Ne.
\]
The exponents \(n\) belong to \(h_e^{-1}\mathbb Z\).
Nearly holomorphic means that each such expansion has
only finitely many negative powers of \(q\).

\begin{theorem}[Cl\'ery--Gritsenko]\label{cyidx:thm:cusps}
Suppose the principal coefficients satisfy
\((h_e/N_e)c_{f/e}(n,\ell)\in\mathbb Z\) whenever
\(4tn-\ell^2\leq0\), at every cusp.
Then
\begin{equation}\label{cyidx:eq:cusp-product}
 q^Ay^Bp^C
 \prod_{f/e}\prod_{(n,m,\ell)>0}
 \left(1-(q^ny^\ell p^{tm})^{N_e}\right)^{
              (h_e/N_e)c_{f/e}(nm,\ell)}
\end{equation}
is a meromorphic modular form for \(\Gamma_t(N)^+\),
with a character or multiplier and weight
\begin{equation}\label{cyidx:eq:cusp-weight}
 k=\frac12\sum_{f/e}\frac{h_e}{N_e}c_{f/e}(0,0).
\end{equation}
Here \(n,m,\ell\in\mathbb Z\), positivity means
\(m>0\), or \(m=0,n>0\),
or \(m=n=0,\ell<0\), and
\[
 A=\frac1{24}\sum_{f/e,\ell}h_ec_{f/e}(0,\ell),\quad
 B=\frac12\sum_{f/e,\ell>0}h_e\ell c_{f/e}(0,\ell),\quad
 C=\frac14\sum_{f/e,\ell}h_e\ell^2c_{f/e}(0,\ell).
\]
The product initially converges in a Weyl chamber.
Its zeros and poles lie on the rational quadratic
divisors prescribed by the negative-discriminant data.
Nonnegative such data give a holomorphic form.
\end{theorem}

This is \cite{GC}*{Theorem~3.1}; the cusp widths in the
integrality hypothesis are part of the statement.
At level two the seed
\(4\vartheta_2(\tau,z)^2/\vartheta_2(\tau,0)^2\)
has constant Laurent row \(y^{-1}+2+y\) at infinity
and constant row \(4\) at zero. The two ratios
\(h_e/N_e\) equal one, so its weight is
\((2+4)/2=3\). The constant term at infinity alone
would give a different weight.

A form with a simple diagonal divisor vanishes precisely
on the translates of \(z=0\), with multiplicity one.
The classification in \cite{GC}*{Theorem~1.4} gives,
up to nonzero scalar, the following eight possibilities.
The final column records the order of the character or
multiplier on \(\Gamma_t(N)\).
\begin{center}
\begin{tabular}{c c c c c}
\(t\)&\(N\)&form&weight&order\\ \hline
1&1&\(\Delta_5\)&5&2\\
2&1&\(\Delta_2\)&2&4\\
3&1&\(\Delta_1\)&1&6\\
4&1&\(\Delta_{1/2}\)&\(1/2\)&8\\
1&2&\(\nabla_3\)&3&2\\
1&3&\(\nabla_2\)&2&2\\
1&4&\(\nabla_{3/2}\)&\(3/2\)&4\\
2&2&\(Q_1\)&1&4
\end{tabular}
\end{center}
The proof first bounds the possibilities by an index
calculation, then constructs these forms by additive and
multiplicative lifts. Its remaining numerical candidate
\((t,N,k)=(2,4,1/2)\) is excluded by its first
Fourier--Jacobi coefficient. Thus a numerical weight
condition alone does not produce an additional form.

\subsection{Divisor-indexed inputs and geometric quotients}

The Mathieu construction has a different input. For an
element \(g\) of \(M_{24}\), let \(\phi_g\) denote its
twisted elliptic genus, of weight zero and index one.
Write the cycle shape of its permutation on twenty-four
points as \(\prod_{r\geq1}r^{b_r}\), so
\(\sum r b_r=24\). Its level \(N_g\) is the congruence
level of the twisted genus; it need not equal the order
of the permutation. The associated product \(\Phi_g\)
uses the whole family
\[
             (\phi_{g^d})_{d\mid N_g}.
\]
The divisor-indexed Jacobi-to-Weil correspondence and
Borcherds lift give \cite{WangWilliamsMathieu}*{Theorem~1.3}
\begin{equation}\label{cyidx:eq:Mathieu-weight}
 \operatorname{wt}(\Phi_g)=\frac12\sum_r b_r-2,
 \qquad L_g=U(N_g)\oplus U\oplus A_1.
\end{equation}
Here \(U\) is the hyperbolic plane, \(U(N_g)\) has
its pairing multiplied by \(N_g\), and \(A_1\)
has a generator of square two. The orthogonal domain
of this lattice of signature \((3,2)\) is the genus-two
Siegel domain. The products are single-valued
meromorphic forms there.

For example, the cycle shapes
\(1^{24},1^82^8,1^63^6,1^42^24^4,1^22^23^26^2\)
give the classes \(1A,2A,3A,4B,6A\) and weights
\(10,6,4,3,2\), respectively. In the full classification
the holomorphic cases are exactly
\(1A,2A,2B,3A,4B\). A meromorphic square root exists
exactly for
\[
 1A,2A,2B,3A,3B,4A,4B,4C,5A,6A,6B,10A,11A.
\]
The additive lift from the first Fourier--Jacobi coefficient
equals \(\Phi_g\) for
\[
                1A,2A,2B,3A,4A,4B,5A,6A,7AB,8A.
\]
For \(3B\) and \(4C\), an additional weight-zero
Borcherds factor occurs \cite{WangWilliamsMathieu}*{Theorem~1.4}.
These assertions concern the stated automorphic forms.
They do not identify every twisted genus with a geometric
automorphism of a chosen K3 surface.

There is a geometric quotient construction when a K3
surface \(S\) carries a symplectic automorphism \(g\)
of order \(N\). Choose an elliptic curve \(E\) and a
point \(a\in E\) of exact order \(N\). Then
\begin{equation}\label{cyidx:eq:CHL}
       X=(S\times E)/\langle(g,t_a)\rangle
\end{equation}
is smooth: every nonidentity power translates \(E\)
by a nonzero point and has no fixed point. The product
volume form descends, so \(K_X\cong\mathcal O_X\).
The elliptic one-form also descends, giving
\(h^{1,0}(X)=1\). In particular the construction uses
a torsion translation, whose order is unrestricted by
the possible orders of linear automorphisms of \(E\).
Symplectic K3 automorphisms of order five supply valid
examples. These quotients and their reduced curve-counting
conjectures are studied in \cite{BryanOberdieckCHL}.
The general primitive counting formula there is a
conjecture; proved base cases retain that narrower scope.

For the untwisted product \(S\times E\), the primitive
Igusa formula is a theorem. To specify its invariants,
take a primitive class \(\beta_h\) on \(S\) with
\(\beta_h^2=2h-2\), choose \(\beta_h^\vee\) pairing
to one with it, and let \(\mathrm{pt}\in H^2(E)\).
Let \(N_{g,h,d}\) be the reduced, disconnected
Gromov--Witten invariant in class \((\beta_h,d)\),
with no contracted connected components and with the
one-point insertion
\(\pi_S^*\beta_h^\vee\cup\pi_E^*\mathrm{pt}\).
In the monic normalization \(\Xi\) of
\eqref{cyidx:eq:D5}, Oberdieck--Pixton prove
\begin{equation}\label{cyidx:eq:OP}
 \sum_{g,h,d\geq0}N_{g,h,d}
       u^{2g-2}q^{h-1}p^{d-1}
       =-\frac1{\Xi(p,q,e^{iu})}.
\end{equation}
The right side is expanded as a Laurent series at
\(u=0\) \cite{OP}*{Theorem~1 and Section~4}.
Their proof shows that Jacobi-form constraints and two
holomorphic-anomaly relations determine the series from
its coefficient \(N_{0,0,0}=1\). The stable-pair
variable is \(-e^{iu}\); this sign matters when comparing
the two counting conventions. Equation~\eqref{cyidx:eq:OP}
is a reduced numerical identity. A Hall product on the
coefficient complexes still requires its correspondences
and compatibility with the reduction.

\section{A second root system: the fake monster}
\label{cyidx:sec:fake}

The fake-monster denominator illustrates the information
carried by a root cone and its imaginary simple roots.
Let \(\Lambda\) be the positive-definite Leech lattice and
write
\[
 II_{25,1}=\Lambda\oplus U,\qquad
 r=(\lambda,m,n),\qquad r^2=\lambda^2-2mn,
 \qquad \rho=(0,0,1).
\]
Its positive roots consist of the vectors of square at most
two with \((r,\rho)<0\), together with \(n\rho\), \(n>0\).
Define \(p_{24}(j)\) and \(\tau(j)\) by
\[
 \eta(\tau)^{-24}=q^{-1}\sum_{j\geq0}p_{24}(j)q^j,
 \qquad \eta(\tau)^{24}=\sum_{j\geq1}\tau(j)q^j.
\]
Set \(p_{24}(j)=0\) for \(j<0\).

The physical-state construction from the lattice vertex
algebra, followed by the quotient by the radical of its
invariant form, gives a generalized Kac--Moody algebra.
For every nonzero root \(r\), its multiplicity is
\(p_{24}(1-r^2/2)\). The real simple roots have
\(r^2=2\), \((r,\rho)=-1\), and each \(n\rho\)
is an imaginary simple root with multiplicity twenty-four
\cite{BorcherdsFake1990}*{Theorems~1 and~3}.
In the completed positive cone its denominator is
\begin{equation}\label{cyidx:eq:fake-denominator}
 \mathbf e^\rho\prod_{r>0}
       (1-\mathbf e^r)^{p_{24}(1-r^2/2)}
  =\sum_{w\in W}\det(w)\sum_{n>0}\tau(n)\mathbf e^{w(n\rho)}.
\end{equation}
The group \(W\) is generated by the real-root reflections.
The imaginary simple roots are mutually orthogonal; their
correction term is the eta product on the right. Thus its
coefficients are those of \(\eta^{24}\), while root
multiplicities come from \(\eta^{-24}\).

The oscillator character also differs from the Leech
lattice-vertex character:
\[
 \eta^{-24}=q^{-1}+24+324q+3200q^2+\cdots,
 \qquad
 \frac{\Theta_\Lambda}{\eta^{24}}
       =j-720=q^{-1}+24+196884q+\cdots.
\]
Here \(\Theta_\Lambda=1+196560q^2+\cdots\).
The coefficient \(196884=196560+324\) displays the
lattice contribution explicitly. Neither character by
itself specifies the Lie bracket used in
\eqref{cyidx:eq:fake-denominator}.

A comparison with the weight-twelve orthogonal Borcherds
form \(\Phi_{12}\) must also preserve automorphic lines.
An ordinary restriction retains weight twelve, so a
nonzero identity with \(D_5^k\) in the same Hodge-line
convention would require \(12=5k\). There is no integral
\(k\) satisfying this. A quasi-pullback is different:
one first divides by the linear equations of the root
divisors containing the subdomain. If there are \(M\)
such factors, and the resulting restriction is defined
and nonzero, its weight is \(12+M\). Its divisor,
character and normalization must then be computed.
The weight obstruction to a raw restriction does not
decide this latter construction.

\section{A convergent determinant model}
\label{cyidx:sec:Fredholm}

The denominator admits an operator realization after a
summability estimate. This construction also makes the
limits of a scalar product identity precise.
For every active positive root \(\alpha=(a,b,\ell)\),
put \(c_\alpha=f(4ab-\ell^2)\). Define a minimal
super vector space \(V_\alpha\) of dimension
\((c_\alpha,0)\) if \(c_\alpha\geq0\), and
\((0,-c_\alpha)\) otherwise. This realizes the
supermultiplicity with total dimension \(|c_\alpha|\).
It need not be the actual root space of
\(\mathfrak n_+\), which may have both parities.

\begin{lemma}\label{cyidx:lem:growth}
For \(n\geq1\), every coefficient of \(q^ny^\ell\)
in \(\phi_{0,1}\) has absolute value at most
\(12\exp(4\pi\sqrt{2n/3})\). Consequently
\begin{equation}\label{cyidx:eq:growth}
 |c_\alpha|\leq12e^{\kappa h(\alpha)},
 \qquad \kappa=2\pi\sqrt{2/3}.
\end{equation}
\end{lemma}
\begin{proof}
Put \(r=|q|^{1/2}\). Apply the Jacobi triple products
to the three theta ratios in \eqref{cyidx:eq:seed}.
The sum of absolute Laurent coefficients of their
expansions is bounded, at \(0<r<1\), by
\[
 12\prod_{j\geq1}\frac{(1+r^j)^4}{(1-r^j)^4}
       \leq12\prod_{j\geq1}(1-r^j)^{-8}.
\]
For \(\vartheta_2\) the factors use the even integers
\(j\) and its initial Laurent polynomial has absolute
coefficient sum one; the other two ratios use the odd
integers. Enlarging each to all \(j\) gives the bound.
For \(r=e^{-t}\),
\[
 \log\prod_{j\geq1}(1-e^{-tj})^{-8}
 =8\sum_{k\geq1}\frac1{k(e^{tk}-1)}
 \leq\frac{4\pi^2}{3t}.
\]
Cauchy's coefficient bound at \(q^n=r^{2n}\)
therefore gives \(12\exp(2nt+4\pi^2/(3t))\).
Minimize at \(t=\pi\sqrt{2/(3n)}\).
For \(n=ab\), the inequality \(2\sqrt{ab}\leq h(\alpha)\)
proves \eqref{cyidx:eq:growth}. The constant row and
the root \((0,0,-1)\) satisfy it directly.
\end{proof}

Give each \(V_\alpha\) an orthonormal homogeneous basis
and form the Hilbert direct sum \(\mathcal H\).
Let \(T_Z\) act on its \(\alpha\)-summand by
\(x_\alpha=p^aq^by^\ell\). For a real \(c>\kappa\)
consider the open chamber in which
\( |x_\alpha|<e^{-c h(\alpha)}\) for every positive root.
It is enough to impose these strict inequalities on the
three real simple roots: every positive charge is their
nonnegative integral combination. It is nonempty: take
\(p=q=e^{-4\pi t}\), \(y=e^{2\pi t}\) with
\(2\pi t>c\). The corresponding imaginary Siegel
matrix is \(t\left(\begin{smallmatrix}2&-1\\-1&2\end{smallmatrix}\right)>0\).

\begin{proposition}\label{cyidx:prop:Fredholm}
On this chamber \(T_Z\) is trace class, locally uniformly
in \(Z\), and
\begin{equation}\label{cyidx:eq:Fredholm}
 \det\nolimits_s(1-T_Z)
 =\frac{\det(1-T_Z|_{\mathcal H_{\bar0}})}
        {\det(1-T_Z|_{\mathcal H_{\bar1}})}
 =\prod_{\alpha>0}(1-x_\alpha)^{c_\alpha}
 =(pqy)^{-1/2}D_5.
\end{equation}
\end{proposition}
\begin{proof}
At fixed height \(h\), there are at most \((h+1)^2\)
root triples: choose \(a,b\leq h\), after which
\(\ell=2(a+b)-h\) is fixed. Thus
\[
 \|T_Z\|_1
 =\sum_\alpha |c_\alpha|\,|x_\alpha|
 \leq12\sum_{h\geq1}(h+1)^2e^{-(c-\kappa)h}<\infty.
\]
On compact subsets the exponent has a positive uniform
margin. The trace-class determinant is therefore
holomorphic and equals the absolutely convergent product
of eigenvalue factors. All these factors are nonzero
in the chamber. Taking their even-to-odd ratio proves
\eqref{cyidx:eq:Fredholm}.
\end{proof}

There is a relative Pfaffian version. In each parity,
set \(D=1-T_Z\) and use the skew operator on
\(\mathcal H\oplus\mathcal H^*\)
\[
 A_Z=\begin{pmatrix}0&D\\-D^{\mathsf T}&0\end{pmatrix}.
\]
Choose the relative orientation for which the Pfaffian
ratio to \(A_0\), where \(D=1\), is \(\det D\).
The finite-dimensional identity follows by pairing each
basis vector with its dual; the common sign cancels in
the ratio. Trace-norm convergence of finite-rank
truncations gives the displayed determinant in the limit.
Taking the ratio between the two parity sectors recovers
\eqref{cyidx:eq:Fredholm}.

This is a bounded identity-plus-trace-class family.
Its skew doubling is not a first-order differential
operator on a manifold. To realize \(D_5\) by such
an operator requires a geometric family, its symbol,
boundary conditions, determinant line and normalization.
Similarly, the modular continuation of the scalar in
\eqref{cyidx:eq:Fredholm} does not supply intertwiners
between the Hilbert spaces in different cusp charts.
Using the actual root superspaces would require a bound
on their total dimensions, in addition to the bound
on their superdimensions proved above.

\section{Elliptic coefficient lines on collision flags}
\label{cyidx:sec:Thom}

We now retain the sheaf from which an elliptic coefficient
map starts. Let \(E\) be a smooth complex elliptic curve,
\(\Theta=\mathcal O_E(0)\), and let \(\vartheta\) be its
divisor section. This \(E\) parametrizes spectral variables.
It is separate from the elliptic factor of a threefold and
from a curve on which fields are inserted.
On \(E^2\), put \(u=z_2-z_1\) and
\(\Delta=\{z_1=z_2\}\).
For two distinct colours joined by two arrows, the elliptic
arrow factor at zero equivariant translations is
\(\vartheta(u)^2\). In the convention that a factor is a
section of the dual kernel line, it gives
\begin{equation}\label{cyidx:eq:pair}
 \mathcal L=u^*\Theta^{-2}=\mathcal O(-2\Delta)
       \xrightarrow{\ \vartheta(u)^2\ }\mathcal O_{E^2}.
\end{equation}
There is no same-colour factor for this dimension vector.
These are the arrow lines in the elliptic shuffle
construction of \cite{YangZhaoElliptic}*{Section~2.1}.

\begin{proposition}\label{cyidx:prop:normal-jet}
The line \(\mathcal L\) is nontrivial. The coefficient
map has image \(\mathcal O(-2\Delta)\) and cokernel
\(\mathcal O_{2\Delta}\). The exact sequence
\begin{equation}\label{cyidx:eq:normal-extension}
 0\longrightarrow\mathcal O_\Delta(-\Delta)
 \longrightarrow\mathcal O_{2\Delta}
 \longrightarrow\mathcal O_\Delta\longrightarrow0
\end{equation}
does not split as a module over the ambient local ring.
\end{proposition}
\begin{proof}
Restrict \(\mathcal L\) to \(\{z_1\}\times E\):
its degree is \(-2\). A formal logarithm at the origin
gives a normal parameter, again denoted \(u\), in which
\(\vartheta(u)=u v(u)\) with \(v\) a unit. Change
the source frame by this unit. The map becomes
\(u^2:\mathbb C[[u]]\to\mathbb C[[u]]\).
Its quotient is \(\mathbb C[[u]]/(u^2)\), and the
submodule \((u)/(u^2)\) gives the first term of
\eqref{cyidx:eq:normal-extension}. A module splitting
would lift \(1\) to an element with nonzero constant
coefficient killed by \(u\), which is impossible.
\end{proof}

The doubled zero remains even when one untwists the
tensor product by changing its coefficient objects.
To see this at every arity, fix a finite simple graph
\(G\) on a set of labels. An edge records an interacting
pair of multiplicity two. For disjoint sets \(A,B\), let
\[
 D_{A,B}=\sum_{\substack{i\in A,j\in B\\\{i,j\}\in E(G)}}
                  \Delta_{ij},\qquad
 \mathcal L_{A,B}=\mathcal O(-2D_{A,B}),
 \qquad
 \mathcal B_A=\mathcal O\left(-2\sum_{
       \substack{i<j,\ i,j\in A\\\{i,j\}\in E(G)}}\Delta_{ij}\right).
\]
All lines are pulled back to the common product of copies
of \(E\). Each has its indicated divisor inclusion.

\begin{proposition}\label{cyidx:prop:flags}
There are canonical, section-preserving isomorphisms
\begin{align}
 \mathcal L_{A,B}\otimes\mathcal L_{A\cup B,C}
 &\cong\mathcal L_{B,C}\otimes\mathcal L_{A,B\cup C},
       \label{cyidx:eq:line-associator}\\
 \mathcal L_{A,B}
 &\cong\mathcal B_{A\cup B}\otimes
             \mathcal B_A^{-1}\otimes\mathcal B_B^{-1}.
       \label{cyidx:eq:line-coboundary}
\end{align}
They satisfy every finite regrouping identity.
The functors \(F\mapsto F\otimes\mathcal B_A^{-1}\)
untwist the labelled tensor product
\((F\boxtimes G)\otimes\mathcal L_{A,B}\).
Replacing every divisor coefficient \(-2\) by \(-1\)
gives compatible square roots of these line identities.
\end{proposition}
\begin{proof}
In \eqref{cyidx:eq:line-associator}, each side contains
exactly once every edge between the three blocks.
Every longer parenthesization counts an edge at the
unique step joining the blocks containing its endpoints.
It therefore gives the same divisor and the same section.
Separating internal edges from cross-edges gives
\eqref{cyidx:eq:line-coboundary}. Tensoring with
\(\mathcal B_{A\cup B}^{-1}\) cancels its first factor
and proves the untwisting assertion. The empty set
has \(\mathcal B_\varnothing=\mathcal O\), so units
are preserved. All the divisor equalities remain true
after halving their coefficients.
\end{proof}

For two labels the transformed coefficient map is
\(\mathcal O\to\mathcal O(2\Delta)\), with the same
doubled zero. Equivalently, multiplication between the
image lines is an isomorphism
\[
 \mathcal B_A\otimes\mathcal B_B\otimes\mathcal L_{A,B}
       \xrightarrow{\sim}\mathcal B_{A\cup B};
\]
the cokernel occurs in the final inclusion
\(\mathcal B_{A\cup B}\hookrightarrow\mathcal O\).
These are labelled coefficient lines. Shuffle
symmetrization, same-colour factors and critical
coefficient maps require the corresponding extra data.

\subsection{Three different derived constructions}

For a commutative ring \(R\) and elements \(f_1,\ldots,f_r\),
write
\[
 K_R(f_1,\ldots,f_r)=
 (R\otimes\Lambda(\epsilon_1,\ldots,\epsilon_r),d),
 \qquad |\epsilon_i|=-1,\qquad d\epsilon_i=f_i.
\]
The differential has degree one and is a graded derivation.
For \(R=\mathbb C[[u]]\), the algebra \(K_R(u^2)\)
has cohomology \(R/(u^2)\) in degree zero. By contrast,
subtracting the two odd generators in \(K_R(u,u)\)
and resolving the remaining regular equation gives
\[
 K_R(u,u)\simeq\mathbb C\otimes\Lambda(\epsilon),
                  \qquad d\epsilon=0.
\]
The first construction has a normal nilpotent; the second
has an excess class in degree minus one.

For a triangle, use formal differences
\(x=z_2-z_1\), \(y=z_3-z_2\). The cone of the total
coefficient product is
\[
 K_{\mathbb C[[x,y]]}\bigl(x^2y^2(x+y)^2\bigr).
\]
It has only degree-zero cohomology, supported on three
doubled lines. The simultaneous derived intersection
of the three pair cokernels is instead
\begin{equation}\label{cyidx:eq:triangle-central}
 K_{\mathbb C[[x,y]]}(x^2,y^2,(x+y)^2).
\end{equation}
It is supported at the origin and has negative cohomology.
It is an incidence object built from the pair maps;
it is not an associativity defect of their tensor product.
Even the ordinary product cone has a filtration, rather
than a tensor-product identity: for nonzero \(f,g\)
in a domain,
\[
 0\longrightarrow R/(f)\xrightarrow{\ g\ }R/(fg)
   \longrightarrow R/(g)\longrightarrow0.
\]
The global version retains the source line twists.

\subsection{A compact coefficient algebra and its trace}

Suppose \(G\) is connected, with \(n\geq1\) vertices
and \(m\) edges. Its spectral translation quotient is
\(B_n=E^n/E_{\mathrm{diag}}\cong E^{n-1}\).
For an edge \(e\), let \(D_e\subset B_n\) be its equality
divisor, let \(L_e=\mathcal O(2D_e)\), and let \(s_e\)
be its divisor section. Define
\begin{equation}\label{cyidx:eq:global-graph}
 \mathscr K_G=\bigotimes_{e\in E(G)}
       [L_e^{-1}\xrightarrow{s_e}\mathcal O_{B_n}],
 \qquad C_G=R\Gamma(B_n,\mathscr K_G).
\end{equation}
The left terms have degree \(-1\). Thus \(\mathscr K_G\)
is a perfect commutative differential graded sheaf.
Its cohomology is supported at the origin: away from
the simultaneous zero locus, an invertible section
contracts its Koszul factor, and connectedness makes
that locus the diagonal. The cohomology sheaves have
finite length and no higher cohomology. Completion at
the origin therefore computes \(C_G\), including its
multiplication.

Adding a constraint gives the exact triangle
\begin{equation}\label{cyidx:eq:edge-triangle}
 \mathscr K_G\otimes L_e^{-1}\xrightarrow{s_e}\mathscr K_G
 \longrightarrow\mathscr K_{G+e}
 \longrightarrow\mathscr K_G\otimes L_e^{-1}[1].
\end{equation}
Locally write an element as \(a+\epsilon_e b\), placing
the new generator first. The last map takes it to \(b\).
The identity \(d(\epsilon_e b)=s_e b-\epsilon_e db\)
gives its shift sign. Inclusion of the even summand is
a map of differential graded algebras. Permuting edge
factors gives the exterior signs and all subset-cube
coherences, including disconnected intermediate graphs.

Choose an oriented spanning tree \(T\). Its
\(k=n-1\) edge differences \(u_1,\ldots,u_k\)
are formal coordinates, and a chord \(c\) has difference
\(\ell_c=\sum_a\sigma_{ca}u_a\), with
\(\sigma_{ca}\in\{-1,0,1\}\). Put \(\beta=m-n+1\).

\begin{theorem}\label{cyidx:thm:finite-graph}
After choosing a formal logarithm and theta frames,
\(C_G\) has the finite model
\begin{equation}\label{cyidx:eq:finite-graph}
 \frac{\mathbb C[u_1,\ldots,u_k]}{(u_1^2,\ldots,u_k^2)}
 \otimes\Lambda(\epsilon_c:c\notin T),\qquad
 d\epsilon_c=2\sum_{a<b}\sigma_{ca}\sigma_{cb}u_au_b.
\end{equation}
Its total dimension is \(2^m\), and its degrees lie in
\([-\beta,0]\). Different trees give quasi-isomorphic
algebras through the full edge Koszul complex.
\end{theorem}
\begin{proof}
The local theta units can be removed by changing the
odd generators. The tree equations form the regular
sequence \(u_1^2,\ldots,u_k^2\). Its Koszul algebra
maps multiplicatively to its quotient. Tensoring with
the free chord factors preserves the quasi-isomorphism.
In the quotient, \(\ell_c^2\) is the differential
in \eqref{cyidx:eq:finite-graph}. Squarefree tree
monomials and exterior chord monomials give
\(2^k2^\beta=2^m\) basis elements. Each tree reduction
starts from the same completed full complex, proving
the comparison without requiring a strict identification
of the finite presentations.
\end{proof}

Fix orders of both sets of generators. The functional
\begin{equation}\label{cyidx:eq:graph-trace}
 \tau(a)=[u_1\cdots u_k\epsilon_1\cdots\epsilon_\beta]a
\end{equation}
gives a perfect pairing \(C_G\otimes C_G\to\mathbb C[\beta]\)
by multiplication. Indeed each basis monomial pairs
with its complementary monomial, with the exterior
permutation sign. The differential lowers exterior
length, so \(\tau(da)=0\). Applying this to \(d(ab)\)
proves the cochain-pairing identity. Hence
\[
 \dim H^{-j}(C_G)=\dim H^{-(\beta-j)}(C_G).
\]
We use \(C[r]^i=C^{i+r}\). A graph symmetry can act
on this normalized trace by a character, which must be
retained before passing to invariant coefficients.
For one vertex the algebra is \(\mathbb C\). The empty
configuration instead has translation quotient \(BE\);
it lies outside this nonempty spectral family.

\section{Displacing the elliptic coefficient}
\label{cyidx:sec:shift}

The doubled factor splits when an equivariant displacement
is retained. Give the two arrows the weight pairs
\((2,0)\) and \((0,2)\). Their compatibility condition
is \(2t_1=2t_2\) in the elliptic group law. On the
formal chart at the origin it forces \(t_1=t_2=t\).
We keep this specialization and put \(h=2t\).
The forward coefficient and its line are
\begin{equation}\label{cyidx:eq:shifted-pair}
 s_h(u)=\vartheta(u)\vartheta(u+h),\qquad
 \mathcal S=u^*\Theta\otimes(u+h)^*\Theta,
 \qquad\mathcal S^{-1}\xrightarrow{s_h}\mathcal O
\end{equation}
on \(E_u\times E_h\). These are the compatible weights
of \cite{YangZhaoElliptic}*{Assumption~2.1 and Remark~2.2}.
They do not produce two independent formal parameters.

The reverse ordering has factor
\(\vartheta(u)\vartheta(u-h)\) in an odd-theta frame.
Write \(D_a=\{u=a\}\). Its image and the forward image
are respectively
\(\mathcal I_-=\mathcal O(-D_0-D_h)\) and
\(\mathcal I_+=\mathcal O(-D_0-D_{-h})\).

\begin{proposition}\label{cyidx:prop:ordered-image}
On a fixed parameter fibre their sum is
\[
 \mathcal J_h=\mathcal I_++\mathcal I_-=
 \begin{cases}
 \mathcal O(-D_0),&2h\ne0,\\
 \mathcal O(-D_0-D_h),&2h=0,\ h\ne0,\\
 \mathcal O(-2D_0),&h=0.
 \end{cases}
\]
In the completed local ring at \((u,h)=(0,0)\), the
relative image ideal is
\begin{equation}\label{cyidx:eq:image-ideal}
 J=(u(u+h),u(u-h))=(u^2,uh).
\end{equation}
The quotient has a nonzero class \([u]\) killed by \(h\).
Formation of the image does not commute with base change
to \(h=0\).
\end{proposition}
\begin{proof}
The sum of divisor ideals on a smooth curve takes the
minimum of their vanishing orders. Comparing the two
divisors gives the three cases. Formal logarithms and
theta frames give the two generators in
\eqref{cyidx:eq:image-ideal}; their sum and difference
give \(u^2,uh\), since two is invertible.
The element \(u\) is absent from \(J\), but \(hu\)
belongs to it. Moreover \(hu\ne0\) in \(J/hJ\):
an equality \(hu=hj\), \(j\in J\), would imply
\(u=j\). Its image in the specialized ambient ring
is zero. Thus \(J/hJ\to\mathbb C[[u]]\) is not
injective, whereas the image of the specialized maps
is the ideal \((u^2)\).
\end{proof}

When \(2h\ne0\), the quotient \(\mathcal J_h/\mathcal I_+\)
is supported at \(-h\), and the reverse quotient at
\(h\). When \(2h=0\), the two image divisors agree.
These statements identify divisor lines and their
inclusions. A full operation also retains its scalar
normalization and equivariant structure.

\subsection{The triangle over the formal parameter disc}

Take three ordered colours with arrows from lower to
higher colour. After removing common translation,
use \(x=z_2-z_1\), \(y=z_3-z_2\). Over
\(R=\mathbb C[[h]]\), the three formal equations are
\begin{equation}\label{cyidx:eq:shifted-equations}
 x(x+h),\qquad y(y+h),\qquad(x+y)(x+y+h).
\end{equation}
Their simultaneous incidence is the Koszul algebra.

\begin{theorem}\label{cyidx:thm:triangle}
It is multiplicatively quasi-isomorphic to
\begin{equation}\label{cyidx:eq:triangle}
 B_h=R[x,y]/(x^2+hx,y^2+hy),\qquad
 A_h=B_h\otimes\Lambda(\eta),\quad |\eta|=-1,
 \quad d\eta=xy.
\end{equation}
Both terms are free of rank four over \(R\).
Its cohomology is free, with bases
\begin{align}
 H^0&=R\{1,x,y\},\qquad x^2=-hx,\quad y^2=-hy,\quad xy=0,
                   \label{cyidx:eq:H0}\\
 H^{-1}&=R\{k_1,k_2,k_3\},\qquad
 k_1=[(x+h)\eta],\quad k_2=[(y+h)\eta],\quad
 k_3=[(xy-h^2)\eta].\label{cyidx:eq:Hminus}
\end{align}
The action on negative classes is
\begin{equation}\label{cyidx:eq:triangle-action}
\begin{array}{c|ccc}
 &k_1&k_2&k_3\\\hline
 x&0&k_3+hk_1&-hk_3-h^2k_1\\
 y&k_3+hk_2&0&-hk_3-h^2k_2.
\end{array}
\end{equation}
Products of two negative classes vanish.
\end{theorem}
\begin{proof}
The first two equations are monic in separate variables
and form a regular sequence. Their Koszul resolution
maps to \(B_h\). In that quotient the third equation
is \(2xy\); divide its odd generator by two. This
gives \eqref{cyidx:eq:triangle} and the multiplicative
quasi-isomorphism. Completion gives the same finite
\(R\)-algebra.

In the basis \(1,x,y,xy\), multiplication by \(xy\)
has only its \(xy\)-row nonzero, equal to
\((1,-h,-h,h^2)\). The cokernel is free on
\(1,x,y\). Its kernel has basis
\(x+h,y+h,xy-h^2\): these elements are killed by
\(xy\), and subtracting their multiples leaves a
scalar killed by the unit entry of the displayed row.
This proves the cohomology assertions. Multiplication
of the representatives proves
\eqref{cyidx:eq:triangle-action}; \(\eta^2=0\)
proves the last assertion.
\end{proof}

The contraction onto cohomology is useful for retaining
the higher operations. Let \(i:H\to A_h\) use the
representatives above. Let \(p:A_h\to H\) fix
\(1,x,y\), kill \(xy,\eta\), and send
\(x\eta,y\eta,xy\eta\) to \(k_1,k_2,k_3\).
Define a degree-minus-one map by
\begin{equation}\label{cyidx:eq:contraction}
 \mathsf H(xy)=\eta,\qquad
 \mathsf H=0\text{ on the other basis elements}.
\end{equation}
Then \(pi=1\), \(d\mathsf H+\mathsf Hd=1-ip\),
and \(\mathsf H^2=\mathsf Hi=p\mathsf H=0\).
For example, \(\mathsf Hd(x\eta)=-h\eta
=x\eta-i(k_1)\), so the parameter term is required.

The functional \(\tau:A_h\to R[1]\) extracting
the coefficient of \(xy\eta\) gives a perfect cyclic
cochain pairing by multiplication. Indeed its matrix
between the two copies of \(B_h\) is
\begin{equation}\label{cyidx:eq:trace-matrix}
 \begin{pmatrix}
 0&0&0&1\\0&0&1&-h\\0&1&0&-h\\1&-h&-h&h^2
 \end{pmatrix},\qquad \det=1.
\end{equation}
The trace kills boundaries, and the derivation rule
gives the cochain identity. Graded commutativity gives
cyclicity. On the displayed cohomology bases the matrix is
\begin{equation}\label{cyidx:eq:H-trace}
 \begin{pmatrix}0&0&1\\0&1&-h\\1&0&-h\end{pmatrix},
                    \qquad\det=-1.
\end{equation}
Thus the induced pairing is perfect over the parameter
ring, including its central fibre.

\subsection{A ternary class killed by the displacement}

Put \(V=Rx\oplus Ry\subset H^0\), and define
\(q:V\otimes V\to R\) by
\(q(x,y)=q(y,x)=1\), \(q(x,x)=q(y,y)=0\).
Let \(\psi(x)=k_1\), \(\psi(y)=k_2\).
We specify the signs of the transferred operations.
Augment both \(A_h\) and \(H\) by sending every
displayed nonunit basis element to zero, and write
\(\overline A\) for the augmentation ideal. On
\(T^c(s\overline A)\), use \(|sa|=|a|-1\) and
\[
 b_1(sa)=-s(da),\qquad
 b_2(sa,sb)=(-1)^{|a|}s(ab),\qquad
 b_3(sa,sb,sc)=(-1)^{2|a|+|b|}s\,m_3(a,b,c).
\]
Each component extends as a degree-one coderivation,
with the Koszul sign from its preceding suspended inputs.

\begin{theorem}\label{cyidx:thm:ternary}
The cohomology algebra extends to a strictly unital
minimal \(A_\infty\) algebra by
\begin{equation}\label{cyidx:eq:m3}
 m_3(a,b,c)=q(b,c)\psi(a)-q(a,b)\psi(c)
                       \quad(a,b,c\in V).
\end{equation}
It is zero on every other basis pattern, and
\(m_n=0\) for \(n\geq4\).
A strictly unital quasi-isomorphism to \(A_h\) has
suspended components
\begin{equation}\label{cyidx:eq:Ainfty-map}
 I_1(sa)=s\,i(a),\qquad
 I_2(sa,sb)=s\mathsf H(i(a)i(b)),\qquad I_{n\geq3}=0.
\end{equation}
\end{theorem}
\begin{proof}
The binary transfer is \(p(i(a)i(b))\). The arity-two
comparison equation is exactly
\(d\mathsf H(i(a)i(b))=i(a)i(b)-i\,m_2(a,b)\).
For three degree-zero inputs, the two binary trees give
\[
 p\bigl(i(a)\mathsf H(i(b)i(c))
              -\mathsf H(i(a)i(b))i(c)\bigr).
\]
The map \(\mathsf H\) extracts the \(xy\) coefficient,
and \(p(x\eta)=k_1\), \(p(y\eta)=k_2\); this is
\eqref{cyidx:eq:m3}. A negative input makes both
terms zero, and unit inputs vanish by the contraction
identities.

In a larger transfer tree, two internal homotopies
either lie on one branch or on different branches.
On one branch the second homotopy encounters \(\eta\)
and vanishes. On different branches the two outputs
multiply to zero by \(\eta^2=0\). Thus all operations
of arity at least four vanish. The final homotopy in
the comparison map also kills an expression containing
\(\eta\), so \(I_{n\geq3}=0\).

In arity three the bar comparison is the displayed
two-tree identity and \(d\mathsf H+\mathsf Hd=1-ip\).
In higher arity its terms contain either two \(I_2\)
outputs or a homotopy on a negative input, and vanish.
Hence the induced coalgebra map satisfies \(b_A I=I b_H\).
Its highest tensor-length component on length \(n\)
is \(i^{\otimes n}\), which is injective. It follows
from \(b_A^2=0\) that \(b_H^2=0\). Finally \(I_1\)
induces the identity on cohomology, proving the
quasi-isomorphism and all arity identities.
\end{proof}

Nonzero formulas can become boundaries after a change
of minimal coordinates. Here the invariant is more
precise. In the normalized Hochschild complex of
\((H,m_2)\), an arity-two cochain of internal degree
\(-1\) is supported on \(V\otimes V\), with values
in \(H^{-1}\). Define
\begin{equation}\label{cyidx:eq:G2}
 G_2(x,x)=0,\quad G_2(x,y)=G_2(y,x)=k_1,
                 \quad G_2(y,y)=k_2.
\end{equation}
On degree-zero inputs its coboundary is
\[
 (\delta G_2)(a,b,c)=aG_2(b,c)-G_2(ab,c)
                       +G_2(a,bc)-G_2(a,b)c.
\]

\begin{proposition}\label{cyidx:prop:ternary-torsion}
The arity-three, internal-degree-minus-one Hochschild
class of \(m_3\) is nonzero and has annihilator exactly
\((h)\). A primitive for its multiple is
\begin{equation}\label{cyidx:eq:hm3}
                      \delta G_2=h\,m_3.
\end{equation}
\end{proposition}
\begin{proof}
The arity-four identity makes \(m_3\) a Hochschild
cocycle. Order the triples as
\[
                 xxx,\ xxy,\ xyx,\ xyy,\ yxx,\ yxy,\ yyx,\ yyy.
\]
Substitution of \eqref{cyidx:eq:triangle-action} gives
\[
 \delta G_2=h(0,k_1,0,-k_2,-k_1,0,k_2,0)=h\,m_3.
\]
Every other normalized pattern vanishes by degree.
Modulo \(h\), the coboundary of any normalized
two-cochain \(G\) at \((x,y,y)\) belongs to
\(\mathbb C k_3\): the products \(xy,y^2\) vanish,
and both \(xH^{-1}\) and \(H^{-1}y\) lie in that
line. But \(m_3(x,y,y)=-k_2\). Hence the class is
nonzero after reduction and therefore over \(R\).
Every scalar outside \((h)\) is a unit; together
with \eqref{cyidx:eq:hm3}, this proves the exact
annihilator.
\end{proof}

At \(h=0\), the same calculation gives the triple
Massey set
\begin{equation}\label{cyidx:eq:Massey}
            \langle[x],[y],[y]\rangle=-k_2+\mathbb C k_3.
\end{equation}
The representative uses \(xy=d\eta\), \(y^2=0\),
and its indeterminacy is \(xH^{-1}+H^{-1}y\).
It excludes zero, so the central algebra is nonformal.
The cohomology modules themselves are free and commute
with base change. The torsion in
Proposition~\ref{cyidx:prop:ternary-torsion} belongs
to the Hochschild class of multiplication.

\subsection{The punctured family and a quadratic comparison}

Over \(K=\mathbb C((h))\), the tree algebra splits into
the four points \(x,y\in\{0,-h\}\). The idempotents are
\begin{equation}\label{cyidx:eq:idempotents}
 \begin{aligned}
 e_{00}&=(x+h)(y+h)/h^2,& e_{10}&=-x(y+h)/h^2,\\
 e_{01}&=-(x+h)y/h^2,& e_{11}&=xy/h^2.
 \end{aligned}
\end{equation}
They are orthogonal and sum to one. On the first three
factors \(xy=0\); on the last \(xy=h^2\).
Projection therefore gives a multiplicative
quasi-isomorphism
\begin{equation}\label{cyidx:eq:generic-triangle}
 A_h\otimes_R K\longrightarrow
       \prod_{(i,j)\ne(1,1)}K\otimes\Lambda(\eta_{ij}),
                    \qquad d=0.
\end{equation}
Multiplication by \(h^{-2}e_{11}\eta\) contracts its
kernel, since its differential is \(e_{11}\).
The poles of the idempotents prevent this decomposition
from extending across the origin.

An Euler connection is explicit on the punctured model:
\begin{equation}\label{cyidx:eq:Euler-connection}
 \nabla_{\partial_h}=\partial_h+
        h^{-1}(x\partial_x+y\partial_y+2\eta\partial_\eta).
\end{equation}
It sends \(x^2+hx\) to \(2h^{-1}(x^2+hx)\), and
similarly for \(y\), so it descends to the quotient.
It commutes with \(d\eta=xy\). The elements
\(x/h,y/h,\eta/h^2\) are horizontal. The residue
weights on \(1,x,y,k_1,k_2,k_3\) are
\(0,1,1,3,3,4\), and analytic continuation around
the puncture has trivial monodromy. This specifies
an algebraic connection; a comparison with geometric
Gauss--Manin transport is an additional construction.

Independent quadratic parameters give a different
deformation. Recentering by \(X=x+h/2\), \(Y=y+h/2\)
gives
\[
 x^2+hx=X^2-h^2/4,\qquad y^2+hy=Y^2-h^2/4,
 \qquad xy=XY-\frac h2(X+Y)+\frac{h^2}{4}.
\]
The linear term in the last expression cannot be dropped.
In fact the obstruction is independent of this choice
of coordinates.

\begin{proposition}\label{cyidx:prop:quadratic-family}
No homomorphism \(\mathbb C[a,b,c]\to R\) makes the
derived pullback of
\[
 \frac{\mathbb C[a,b,c][X,Y]}{(X^2-a,Y^2-b)}
           \otimes\Lambda(\zeta),\qquad |\zeta|=-1,
           \qquad d\zeta=XY-c,
\]
quasi-isomorphic to \(A_h\) as a differential graded
\(R\)-algebra.
\end{proposition}
\begin{proof}
The displayed family is termwise free over the parameter
ring, so ordinary tensor product computes its derived
pullback. Over \(K\), multiplication by \(XY-c\) has
matrix
\begin{equation}\label{cyidx:eq:quadratic-matrix}
 F=\begin{pmatrix}
 -c&0&0&ab\\0&-c&b&0\\0&a&-c&0\\1&0&0&-c
 \end{pmatrix}
\end{equation}
in the basis \(1,X,Y,XY\). To have three-dimensional
kernel and cokernel it must have rank one. Its
two-by-two minors include \(a,b,c\), up to signs,
so all three must vanish in \(K\), hence in \(R\).
The pullback would be the constant central algebra.
Its Massey set \(-k_2+Kk_3\) still excludes zero.
But \eqref{cyidx:eq:generic-triangle} makes the generic
fibre of \(A_h\) formal. This contradiction proves
the assertion for every base map.
\end{proof}

\section{The whole spectral parameter curve}
\label{cyidx:sec:global-shift}

The formal family comes from a global perfect complex.
On \(E_x\times E_y\times E_h\), set
\(u_1=x,u_2=y,u_3=x+y\),
\(\mathcal S_i=u_i^*\Theta\otimes(u_i+h)^*\Theta\),
and \(s_i=\vartheta(u_i)\vartheta(u_i+h)\).
For the projection \(\pi\) to \(E_h\), define
\begin{equation}\label{cyidx:eq:global-shift}
 \mathscr K=\bigotimes_{i=1}^3
      [\mathcal S_i^{-1}\xrightarrow{s_i}\mathcal O],
                      \qquad \mathscr C=R\pi_*\mathscr K.
\end{equation}
Write \(E[n]=\ker([n]:E\to E)\).

\begin{theorem}\label{cyidx:thm:global-fibres}
The complex \(\mathscr C\) has two locally free terms
of rank four on \(E_h\), in degrees \(-1,0\).
Its derived fibres have dimensions
\[
\begin{array}{c|cc}
 h&\dim H^0&\dim H^{-1}\\\hline
 0&3&3\\
 h\ne0,\ 2h\ne0&3&3\\
 h\ne0,\ 2h=0&4&4.
\end{array}
\]
Near a nonzero two-torsion point, three summands have
zero differential. The fourth is
\([\mathcal O\xrightarrow{v}\mathcal O]\), where
\(v\) is a local parameter times a unit. It contributes
a skyscraper to \(H^0(\mathscr C)\); the additional
negative class of its derived fibre is a
\(\operatorname{Tor}_1\) class.
\end{theorem}
\begin{proof}
The first two sections cut out
\[
 Z=\{x=0\text{ or }-h\}\times_{E_h}
                 \{y=0\text{ or }-h\},
\]
with their Cartier multiplicities. Each factor is
finite flat of degree two: it is the sum of two
sections, with a monic quadratic equation near their
meeting point and no vertical component. Their
product is finite flat of degree four, and the first
two equations form a regular sequence. The complex
therefore reduces to the finite pushforward of
\([\mathcal S_3^{-1}|_Z\xrightarrow{s_3}\mathcal O_Z]\).
An invertible sheaf on this finite flat scheme has
locally free pushforward of rank four. The resulting
two-term complex computes every derived fibre.

For \(h\ne0\), the four points of \(Z_h\) are
distinct. The third section vanishes at
\((0,0),(0,-h),(-h,0)\). At \((-h,-h)\) its value
is, in the relevant line,
\[
                    \vartheta(-2h)\vartheta(-h).
\]
This is nonzero unless \(2h=0\). At a nonzero
two-torsion point the second factor is nonzero, and
the first has a simple zero because multiplication
by two is étale. Local frames give the
stated fourth complex. Its kernel over the base is
zero and its cokernel is \(\mathcal O/(v)\).
Derived specialization supplies the additional
degree-minus-one class. At \(h=0\), formal logarithms
reduce to \eqref{cyidx:eq:triangle}, whose cohomology
has ranks three and three.
\end{proof}

Every nonzero fibre is formal: each surviving reduced
tree point carries one exterior generator with zero
differential, and a nonsurviving point is contractible.
At the three nonzero points of \(E[2]\), four exterior
factors survive. This is a local description and does
not choose global frames for the persistent summands.
On the original parameter cover \(h=2t\), the extra
fibres occur at the twelve points
\(E[4]\setminus E[2]\); the nonformal doubled fibre
occurs at the four points of \(E[2]\).

\subsection{Square roots on the halving cover}

Rigidify the symmetric line \(\Theta\) at zero.
The coefficient line
\(\mathcal S=u^*\Theta\otimes(u+h)^*\Theta\)
has degree two on each \(u\)-fibre, and hence four
degree-one square roots on that fibre. These roots
do not vary continuously over the original base.

\begin{theorem}\label{cyidx:thm:halving-cover}
There is no line \(\mathcal R\) on \(E_u\times E_h\)
with \(\mathcal R^2\cong\mathcal S\), nor a continuous
relative choice of degree-one root, even after a base
line is tensored with \(\mathcal S\).
On the cover \(h=2t\), the cube theorem gives
\begin{equation}\label{cyidx:eq:parallelogram}
 u^*\Theta\otimes(u+2t)^*\Theta
       \cong(u+t)^*\Theta^2\otimes t^*\Theta^2.
\end{equation}
Thus
\begin{equation}\label{cyidx:eq:root-line}
                  \mathcal R_t=(u+t)^*\Theta\otimes t^*\Theta
\end{equation}
is a root line. A nonzero deck shift \(t\mapsto t+\xi\),
\(\xi\in E[2]\), changes its restriction to a \(u\)-fibre
by the corresponding nontrivial two-torsion line.
For \(2t\ne0\), the section
\(\vartheta(u)\vartheta(u+2t)\) is not the square
of a regular section of \(\mathcal R_t\).
\end{theorem}
\begin{proof}
Restrict \(\mathcal S\) to \(\{u_0\}\times E_h\).
Its degree is one, whereas the square of a line has
even degree. This rules out an absolute root.
Under \(\operatorname{Pic}^1(E)\cong E\), a root
on a fixed fibre is a point \(m\) with \(2m=-h\).
A continuous choice would induce
\(2m_*=-\mathrm{id}\) on \(\pi_1(E)=\mathbb Z^2\),
which is impossible. Tensoring by a base line does
not change these relative Picard classes.

For a symmetric rigidified line the cube theorem
implies
\[
 (f+g)^*\Theta\otimes(f-g)^*\Theta
       \cong(f^*\Theta)^2\otimes(g^*\Theta)^2.
\]
The rigidification removes the constant fibre factor;
see \cite{StacksCube}. Apply this identity to
\(f=u+t\), \(g=t\), proving
\eqref{cyidx:eq:parallelogram} and its specified root.
The ratio after a deck shift restricts to a translate
of \(T_\xi^*\Theta\otimes\Theta^{-1}\). The
degree-one polarization identifies \(E\) with
\(\operatorname{Pic}^0(E)\), so this is nontrivial
of order two for \(\xi\ne0\).
Finally the two zeros \(0,-2t\) of the section are
distinct and simple. A square of a regular section
has even vanishing order at every point.
\end{proof}

Equation~\eqref{cyidx:eq:parallelogram} is an identity
of lines. It is not the scalar identity
\(\vartheta(u)\vartheta(u+2t)
  =(\vartheta(u+t)\vartheta(t))^2\), whose right
side vanishes identically at \(t=0\).
For an ordered flag, tensor the root
\eqref{cyidx:eq:root-line} over its interacting
cross-pairs, using the fixed square isomorphism for
each labelled pair. Every parenthesization uses each
pair once. Associativity of tensor products therefore
gives all flag coherences, and squaring gives the
coefficient-line identifications. Negative kernel
lines use the inverse roots. Same-colour and dynamical
factors require their own line constructions.

\section{Shifted graph coefficients and integral torsion}
\label{cyidx:sec:shifted-graphs}

Let \(G\) be connected on ordered vertices
\(0,\ldots,n-1\), and orient every edge from lower
to higher vertex. Set \(z_0=0\) and assign the edge
\(ij\) the equation
\((z_j-z_i)(z_j-z_i+h)\).
Choose an oriented spanning tree, with coordinates
\(u_1,\ldots,u_k\), \(k=n-1\), and retain the
signed chord paths \(\ell_c=\sum_a\sigma_{ca}u_a\).
Let \(\beta=|E(G)|-n+1\).

\begin{theorem}\label{cyidx:thm:shifted-graphs}
The full edge Koszul algebra is multiplicatively
quasi-isomorphic to
\begin{equation}\label{cyidx:eq:shifted-graph-model}
 C_{G,h}=\frac{R[u_1,\ldots,u_k]}
                     {(u_a(u_a+h):1\leq a\leq k)}
                 \otimes\Lambda(\epsilon_c:c\notin T),
\end{equation}
where \(|\epsilon_c|=-1\) and
\begin{equation}\label{cyidx:eq:shifted-chord}
 d\epsilon_c=2\sum_{a<b}\sigma_{ca}\sigma_{cb}u_au_b
                  +h\sum_a(\sigma_{ca}-\sigma_{ca}^2)u_a.
\end{equation}
Its degree \(-j\) term is free of rank
\(2^k\binom\beta j\), and its total rank is
\(2^{|E(G)|}\). Define
\[
 \mathcal A_G=\{s\in\{0,1\}^k:
                \textstyle\sum_a\sigma_{ca}s_a\in\{0,1\}
                        \text{ for every chord }c\}.
\]
After inverting \(h\), projection gives a
multiplicative quasi-isomorphism
\begin{equation}\label{cyidx:eq:graph-generic}
 C_{G,h}\otimes_R K\longrightarrow
      \prod_{s\in\mathcal A_G}
             \Lambda_K(\epsilon_1,\ldots,\epsilon_\beta),
                         \qquad d=0.
\end{equation}
\end{theorem}
\begin{proof}
Tree differences are integral unimodular coordinates
on the translation quotient. Their monic quadratics
form a regular sequence; quotienting by them gives
the multiplicative reduction as in
Theorem~\ref{cyidx:thm:finite-graph}. Expanding
\(\ell_c^2+h\ell_c\) and using \(u_a^2=-hu_a\)
gives \eqref{cyidx:eq:shifted-chord}. In particular
a negatively traversed edge contributes \(-2hu_a\).
Squarefree tree monomials and exterior subsets give
the ranks.

Over \(K\), the tree algebra splits into points
\(u_a=-hs_a\). On the \(s\)-factor a chord
differential becomes \(h^2L_c(s)(L_c(s)-1)\),
where \(L_c(s)=\sum_a\sigma_{ca}s_a\).
If one is nonzero, it is invertible and multiplication
by its reciprocal times \(\epsilon_c\) contracts
that factor. On every remaining factor all
differentials are zero, proving the projection formula.
\end{proof}

The central and generic dimensions do not determine
the order of integral torsion. For the following
graphs it can be computed from their finite matrices.
Write \(ij\) for the edge \(\{i,j\}\); dimensions
are listed in degrees \(0,-1,\ldots\).

\begin{proposition}\label{cyidx:prop:graph-torsion}
The shifted graph algebras have the following cohomology.
\begin{center}
\begin{tabular}{c c c c}
\(G\)&central dimensions&generic dimensions&\(R\)-torsion\\\hline
\(K_3\)&\((3,3)\)&\((3,3)\)&0\\
\(01,12,23,03\)&\((6,6)\)&\((4,4)\)&\((R/(h))^2\) in \(H^0\)\\
\(K_4\setminus\{01\}\)&\((5,10,5)\)&\((5,10,5)\)&0\\
\(K_4\)&\((4,14,14,4)\)&\((4,12,12,4)\)&\((R/(h))^2\) in \(H^{-1}\)
\end{tabular}
\end{center}
For the third graph one may give the vertices colours
\((1,1,2,3)\) and include only intercolour factors.
The last graph has four distinct colours.
\end{proposition}
\begin{proof}
Use the ordered trees and chord paths
\[
\begin{array}{c|c|c}
G&T&(\sigma_c)_c\\\hline
K_3&(01,12)&(1,1)\\
01,12,23,03&(01,03,23)&(-1,1,-1)\\
K_4\setminus\{01\}&(02,12,13)&(1,-1,1),\ (0,-1,1)\\
K_4&(01,12,23)&(1,1,0),\ (1,1,1),\ (0,1,1).
\end{array}
\]
There are respectively \(3,4,5,4\) admissible bit
assignments. Theorem~\ref{cyidx:thm:shifted-graphs}
gives the generic dimensions.

For the integral calculation, order squarefree
monomials by bit mask, with \(u_1\) the least
significant bit, and exterior subsets lexicographically.
Multiplication by a tree coordinate has the matrix
\(\left(\begin{smallmatrix}0&0\\1&-h\end{smallmatrix}\right)\)
on that coordinate and the identity on the others.
Equation~\eqref{cyidx:eq:shifted-chord} then specifies
every matrix. Removing the \(r\)-th exterior factor
has sign \((-1)^{r-1}\).

Cancel a nonzero constant entry \(w=D_{rc}\) by
replacing the remaining entries with
\(D_{ij}-D_{ic}w^{-1}D_{rj}\). Apply the inverse
basis changes to adjacent differentials. Their
components in the cancelled summand vanish because
\(d^2=0\), so this splits off
\([R\xrightarrow{1}R]\). Reading from lowest
cochain degree upwards and then by rows, the first
available constant entries give respectively
\(1,2,6,14\) cancellations. For the first and third
graphs the remaining differential is zero, with
ranks \((3,3)\) and \((5,10,5)\).

For the four-cycle the remaining map is \(6\times6\).
After factoring out \(h\), its nonzero rows are
\[
 \begin{pmatrix}
 -2h&2h^2&4h&-4h^2&2h^2&-2h^3\\
 -2&2h&-2&2h&2h&-2h^2\\
 -2&2h&4&-4h&2h&-2h^2
 \end{pmatrix}.
\]
The first row is \(h\) times the third; the last
two have a constant invertible minor in columns
one and three. Every other column is in their
\(R\)-span. The two nonzero Smith factors of the
original residual map are therefore \(h,h\).

For \(K_4\), the remaining ranks are
\((4,14,14,4)\). Only the map from degree \(-2\)
to degree \(-1\) is nonzero. Dividing it by \(h\)
and removing zero rows and columns gives
\[
 \begin{pmatrix}
 -2&2h&2&0&-2h&-2h&0\\
 2&-2h&0&-2&2h&0&2h\\
 0&0&-2&2&0&2h&-2h
 \end{pmatrix}.
\]
The rows sum to zero. The first two have an invertible
constant minor in columns one and three, and all
remaining columns are in their span. Its nonzero
Smith factors are again \(h,h\). Thus the two
torsion summands now lie in degree \(-1\).
Every residual differential vanishes modulo \(h\),
which gives the central dimensions and excludes
higher torsion orders.
\end{proof}

The repeated-colour example concerns only its indicated
intercolour factors; a Hall coefficient must also
include the same-colour contribution. The number of
admissible bit assignments is a count of branches
of this spectral incidence algebra. Identifying it
with a compact BPS multiplicity or an Igusa coefficient
requires a map from the corresponding compact theory.

\section{The compact comparison and the weight-five line}
\label{cyidx:sec:comparison}

The constructions in this chapter retain several
different kinds of coefficients. The Hilbert-scheme
elliptic complexes yield \(F_T\) and \(F_S\) after
taking Euler supercharacters. The root Lie algebra
has a bracket, and its Chevalley differential gives
the inverse enveloping character. The spectral
incidence algebra has its cyclic trace, normal jets
and the relative ternary class \([m_3]\). The
automorphic form is a section of
\(\lambda^5\otimes\nu_5\).

For a K3 surface \(S\) and an elliptic curve
\(E_{\mathrm{geom}}\), a reduced oriented critical
Hall theory on \(S\times E_{\mathrm{geom}}\) must
be defined on its compact moduli, with extension
correspondences and their common translation torsors.
The empty sector lies over \(BE_{\mathrm{geom}}\).
Its translation reduction differs from the quotient
of spectral coordinates used in
\eqref{cyidx:eq:global-graph}. The dual-number-valued
reduced integration of \cite{OberdieckShen} and the
local trace \eqref{cyidx:eq:trace-matrix} have distinct
domains and coefficient rings.

\begin{problem}\label{cyidx:prob:compact}
Construct a comparison from independently defined
reduced oriented compact critical coefficients to
a collision recipient carrying
\eqref{cyidx:eq:global-shift} over the whole spectral
parameter curve. Determine whether a multiplicative
comparison identifies the two ordered coefficient
lines, their exceptional fibres, and the halving
cover with its \(E[2]\)-action on root choices.
A proposed equivalence of coefficient algebras must
identify the minimal-model class \([m_3]\) and its
integral relation \(\delta G_2=h\,m_3\).
\end{problem}

Replacing \(G_2\) by \(G_2/h\) changes the coefficient
ring and cannot give a regular comparison over the
whole disc. A representation on a state complex has
another target: its equations hold in the endomorphism
algebra of that state and need not preserve the
auxiliary incidence multiplication.

\begin{problem}\label{cyidx:prob:weight-five}
Construct a geometric family over genus-two moduli
whose determinant or Pfaffian line is
\(\lambda^5\otimes\nu_5\), and identify its section
with the normalized \(D_5\). Relate its degeneration
and sewing maps to the compact critical coefficients
and their elliptic collision lines.
\end{problem}

The line \eqref{cyidx:eq:root-line} on a fixed spectral
elliptic curve supplies a square root with explicit
deck action. It does not by itself give a family over
genus-two moduli. The determinant model
\eqref{cyidx:eq:Fredholm} supplies a convergent scalar
in a cusp chamber. A solution of the last problem
must additionally construct the geometric family,
its transition maps and the operation-preserving
comparisons.
